% Propietario conservado: /Users/ruben/Documents/New project/output/REVISION_ARGUMENTAL_APERTURAS_CIERRES_20260915/VII/delivery/MANUSCRIPT_VII_ES_EN_COMPLETE/source_es/manuscrito/30e_corriente_extension_minima.tex
% Artículo VII. Incorporación aditiva, 10-09-2026.
% RESULTADO_RECUPERADO: fuentes_conservadas/18_clausura_energetica.tex,
% líneas 357--485: Gramiano, peso, extensión mínima y eliminación de Schur;
% líneas 149--180 y 235--295: normalización y unicidad de la respuesta.
% Residencia anterior en VII: 30_pantalla_y_respuesta.tex,
% vii:prop:incidencia, vii:eq:lambda-pantalla, vii:prop:minima-extension.
% FORMALIZACION_NUEVA / CERTIFICADO_NUEVO: diferencial del mínimo con
% incidencia, frontera y normalización variables; incidencia transportada
% por bloques y composición con coordenadas de conexión.
% La familia de rutas que realiza cada incidencia se mantiene explícita.
% No se atribuyen al propietario histórico ni esa familia diferenciable
% completa ni la evaluación de una amplitud material de espín.
% Dependencias: 30, 30c y la convención variacional de 31. No modifica
% esos fragmentos ni constituye por sí solo una realización material.

\section{Diferencial de la extensión mínima de frontera}
\label{vii:vii:sec:corriente-extension-minima}

La incidencia cúbica determina simultáneamente el campo de extensión
mínima y el peso de su respuesta. Para un dato de frontera, ambos se
calculan mediante el inverso del Gramiano de incidencia. Se desarrolla
ahora el diferencial de esa misma operación, conservando las variaciones
de la frontera, del transporte y de la escala de normalización. De este
modo, el campo que interviene en la respuesta queda seleccionado por
el problema variacional de pantalla.

\subsection{Campo minimizante y peso de respuesta}
\label{vii:vii:subsec:extension-minima-parametrizada}

Sean \(H\) y \(Q\) espacios hermíticos de dimensión finita, con
productos fijos y antilineales en la primera variable. El caso real se
obtiene sustituyendo los adjuntos por transpuestas. En un abierto de
parámetros reales \(\Omega\), consideremos familias de clase \(C^1\)
\[
 \mathsf C(\theta):H\longrightarrow Q,\qquad
 b(\theta)\in Q,\qquad \chi(\theta)>0,
 \qquad \operatorname{Ran}\mathsf C(\theta)=Q.
\]
La sobreyectividad equivale a rango fila completo en bases ortonormales.
Definimos el Gramiano, el potencial conjugado y la extensión
\begin{equation}
 \mathsf L=\mathsf C\mathsf C^*,\qquad
 y=\mathsf L^{-1}b,\qquad f_{\min}=\mathsf C^*y,
 \qquad \Lambda=\chi\mathsf L^{-1}.
 \label{vii:vii:eq:gramiano-extension-minima}
\end{equation}
La incidencia \(\mathsf C\) y el emparejamiento celular \(M\) de
\eqref{vii:vii:eq:coordenadas-corriente-material} son operadores diferentes.

\begin{proposition}[Selección variacional de la extensión]
\label{vii:vii:prop:minimo-parametrizado}
La extensión \(f_{\min}\) es la única minimizante de la norma entre
los campos \(f\in H\) que satisfacen \(\mathsf C f=b\). Su energía es
\begin{equation}
 E(\theta)=\chi\min_{\mathsf C f=b}\|f\|^2
 =\chi\langle b,\mathsf L^{-1}b\rangle
 =\langle b,\Lambda b\rangle.
 \label{vii:vii:eq:energia-extension-minima}
\end{equation}
El campo \(f_{\min}\), el potencial \(y\) y la energía \(E\)
dependen de manera \(C^1\) de los parámetros.
\end{proposition}

\begin{proof}
La sobreyectividad de \(\mathsf C\) hace inyectivo a \(\mathsf C^*\),
por lo que \(\langle z,\mathsf Lz\rangle=\|\mathsf C^*z\|^2>0\)
para \(z\ne0\). Así, \(\mathsf L\) es invertible y
\(\mathsf C f_{\min}=\mathsf L y=b\). Toda solución es
\(f_{\min}+z\), con \(z\in\ker\mathsf C\). Como
\(f_{\min}\in\operatorname{Ran}\mathsf C^*=(\ker\mathsf C)^\perp\),
\[
 \|f_{\min}+z\|^2=\|f_{\min}\|^2+\|z\|^2,
 \qquad
 \|f_{\min}\|^2=\langle y,\mathsf L y\rangle
 =\langle b,\mathsf L^{-1}b\rangle.
\]
Estas identidades prueban el mínimo, su unicidad y su valor. La inversión
es una operación diferenciable en el abierto de matrices invertibles;
las expresiones de \eqref{vii:vii:eq:gramiano-extension-minima} proporcionan
la regularidad afirmada.
\end{proof}

La realización cúbica de \ref{vii:vii:prop:incidencia} corresponde a
\(H=\mathbb R^{\{-1,1\}^3}\), al cociente de pantalla
\(Q=P_\square\mathbb R^6\) y a la incidencia \(\mathsf C=M:H\to Q\).
Su Gramiano y su normalización son
\begin{equation}
 \mathsf L_\square=12P_u+4P_-,\qquad
 \chi_m=\frac{112}{81\ell_m^2},\qquad
 \Lambda_m=\frac{28}{243\ell_m^2}(P_u+3P_-).
 \label{vii:vii:eq:normalizacion-minimo-cubico}
\end{equation}
El factor \(112/81\) procede de la soldadura y del defecto cuadrático
de compresión de \eqref{vii:vii:eq:lambda-pantalla}. Si \(b\) tiene dimensión
de longitud, también la tienen \(y\) y \(f_{\min}\), la incidencia es
adimensional y \(E\) es adimensional. La conversión a acción utiliza
la normalización en la recta de acción que se especifica más adelante.

\subsection{Diferencial completo del mínimo}
\label{vii:vii:subsec:diferencial-minimo}

\begin{theorem}[Respuesta a la variación de incidencia y frontera]
\label{vii:vii:thm:diferencial-minimo}
Para toda dirección tangente real \(\delta\theta\), se tiene
\begin{equation}
 \delta E
 =\delta\chi\,\langle b,y\rangle
   +2\chi\operatorname{Re}
       \langle y,\delta b-(\delta\mathsf C)f_{\min}\rangle.
 \label{vii:vii:eq:diferencial-minimo-completo}
\end{equation}
Si \(\chi=112/(81\ell_m^2)\), con \(\ell_m>0\) variable,
\begin{equation}
 \delta E=-2E\frac{\delta\ell_m}{\ell_m}
 +2\chi_m\operatorname{Re}
       \langle y,\delta b-(\delta\mathsf C)f_{\min}\rangle.
 \label{vii:vii:eq:diferencial-minimo-escala}
\end{equation}
\end{theorem}

\begin{proof}
Diferenciar \(\mathsf L\mathsf L^{-1}=I_Q\) da
\[
 \delta\mathsf L^{-1}
 =-\mathsf L^{-1}(\delta\mathsf L)\mathsf L^{-1},\qquad
 \delta\mathsf L=(\delta\mathsf C)\mathsf C^*
                      +\mathsf C(\delta\mathsf C)^*.
\]
Aplicando la regla del producto a \eqref{vii:vii:eq:energia-extension-minima}
y usando la autoadjunción de \(\mathsf L^{-1}\), resulta
\[
 \delta E=\delta\chi\,\langle b,y\rangle
 +2\chi\operatorname{Re}\langle y,\delta b\rangle
 -\chi\langle y,(\delta\mathsf L)y\rangle.
\]
Los dos términos del último producto son conjugados entre sí y
\(\mathsf C^*y=f_{\min}\); por tanto,
\[
 \langle y,(\delta\mathsf L)y\rangle
 =2\operatorname{Re}\langle y,(\delta\mathsf C)f_{\min}\rangle.
\]
Se obtiene \eqref{vii:vii:eq:diferencial-minimo-completo}.
Finalmente, \(\delta\chi_m=-2\chi_m\delta\ell_m/\ell_m\)
produce \eqref{vii:vii:eq:diferencial-minimo-escala}.
\end{proof}

La fórmula incorpora la variación del campo minimizante mediante la
ecuación que lo determina. También se comprueba directamente:
\(\mathsf C\delta f_{\min}=\delta b-(\delta\mathsf C)f_{\min}\)
y \(f_{\min}=\mathsf C^*y\) implican
\[
 \delta\|f_{\min}\|^2
 =2\operatorname{Re}\langle f_{\min},\delta f_{\min}\rangle
 =2\operatorname{Re}
     \langle y,\delta b-(\delta\mathsf C)f_{\min}\rangle.
\]
Por ejemplo, si \(\mathsf C\) es fijo y \(b=\ell_m b_0\), con
\(b_0\) fijo, los dos términos de
\eqref{vii:vii:eq:diferencial-minimo-escala} se cancelan y la energía
permanece constante. Este control conserva conjuntamente la dimensión
de frontera y su normalización.

\subsection{Cociente fijo y cambios de marco}
\label{vii:vii:subsec:marcos-minimo}

Las derivadas anteriores están tomadas en espacios \(H,Q\) y productos
hermíticos fijos. En la presentación de seis caras, la inversa siempre
actúa sobre \(Q=P_\square H_F\). Para una familia de cocientes
\(Q(\theta)\), se fija primero una trivialización isométrica
\(V_Q(\theta):Q\to Q(\theta)\), y análogamente para el dominio.
Las derivadas de esos mapas se incluyen al diferenciar las expresiones
transportadas a los espacios fijos. Las fórmulas se aplican en cada
abierto donde se conserve el rango. Los cambios del producto energético
requieren diferenciar además ese producto; una métrica lorentziana
y el producto positivo usado para minimizar conservan aquí sus tipos
distintos.

\begin{proposition}[Covariancia del mínimo]
\label{vii:vii:prop:covariancia-minimo}
Sean \(g_H,g_Q\) unitarios. Bajo
\[
 \mathsf C'=g_Q\mathsf C g_H^*,\qquad b'=g_Qb,\qquad \chi'=\chi,
\]
se cumple \(y'=g_Qy\), \(f'_{\min}=g_Hf_{\min}\) y \(E'=E\).
Cuando los cambios de marco dependen del parámetro, el diferencial
completo conserva esa igualdad. En particular, una variación pura de
marco tiene \(\delta E=0\).
\end{proposition}

\begin{proof}
El Gramiano transformado es \(\mathsf L'=g_Q\mathsf Lg_Q^*\),
de donde se obtienen las transformaciones de \(y,f_{\min}\) y
la igualdad de energías. Para una variación pura de marco, en un punto
con marcos identidad, sean \(A_Q^*=-A_Q\) y \(A_H^*=-A_H\)
sus generadores. Entonces
\[
 \delta\mathsf C=A_Q\mathsf C-\mathsf C A_H,\qquad
 \delta b=A_Qb,\qquad
 \delta b-(\delta\mathsf C)f_{\min}=\mathsf C A_Hf_{\min}.
\]
Su contribución a \eqref{vii:vii:eq:diferencial-minimo-completo} es
\(2\chi\operatorname{Re}\langle f_{\min},A_Hf_{\min}\rangle=0\).
La prueba incluye, así, los términos debidos a los marcos móviles.
\end{proof}

\subsection{Incidencia transportada y componentes de conexión}
\label{vii:vii:subsec:incidencia-conexion-minima}

La dependencia de conexión se especifica antes de efectuar la
contracción. Una realización explícita de la incidencia por bloques
se obtiene como sigue. Sean \(\mathcal K\) una fibra hermítica común,
\(F=\{(i,\epsilon):1\leq i\leq3,\ \epsilon=\pm1\}\) y
\(V=\{-1,1\}^3\). Para cada incidencia \(\nu_i=\epsilon\),
la realización utilizada declara una ruta \(\gamma_{i,\epsilon,\nu}\)
desde la fibra del vértice hasta la de la cara, con transporte
\(U_{i,\epsilon,\nu}\). Las trivializaciones identifican esas fibras
con \(\mathcal K\). Definimos
\begin{equation}
 \begin{split}
 (\mathsf B_Uf)_{i,\epsilon}
   &=\sum_{\nu_i=\epsilon}U_{i,\epsilon,\nu}f_\nu,\\
 \mathsf C_U&=(q_\square\otimes I_{\mathcal K})\mathsf B_U:
   \bigoplus_{\nu\in V}\mathcal K\longrightarrow Q\otimes\mathcal K,
 \qquad q_\square=P_\square:H_F\to Q.
 \end{split}
 \label{vii:vii:eq:incidencia-transportada-minima}
\end{equation}
Esta es una continuación variacional explícita de la incidencia cúbica
en la representación de rutas declarada. Para transportes identidad,
\(\mathsf C_U=M\otimes I_{\mathcal K}\) y se recupera
\eqref{vii:vii:eq:normalizacion-minimo-cubico}, tensorizada con la fibra.
El Gramiano sigue siendo positivo en un entorno de esa realización,
pues la positividad definida es una condición abierta. En cada dominio
de rango completo, el peso y el campo se fijan por
\eqref{vii:vii:eq:gramiano-extension-minima}. Las rutas que realizan las
incidencias, sus orientaciones y sus memorias son datos de esa
realización del transporte APP--TRIT--TPK y acompañan a la construcción.

Si \(\omega=(\omega_j)\) son coordenadas reales de conexión,
la derivada de \(\mathsf C_U\) se calcula sustituyendo en cada bloque
\begin{equation}
 \partial_j U_\gamma
 =\sum_{k=1}^{N}
    U_N\cdots U_{k+1}(\partial_jU_k)U_{k-1}\cdots U_1,
 \qquad U_\gamma=U_N\cdots U_1.
 \label{vii:vii:eq:derivada-incidencia-rutas}
\end{equation}
Los productos vacíos son identidades. La fórmula procede de la regla
del producto y conserva el lugar de cada operación en la ruta.
Así, \(\mathsf C_j:=\partial_j\mathsf C_U\) se obtiene de una
suma finita de matrices conocidas en la realización elegida, seguida
de la proyección \(q_\square\otimes I\). Si se utiliza un cociente
móvil, su derivada se incorpora conforme a
\ref{vii:vii:subsec:marcos-minimo}. La matriz cúbica entera fija corresponde
al caso \(\mathsf C_j=0\); la respuesta a conexión procede de la
dependencia de los transportes efectivamente declarada en
\eqref{vii:vii:eq:incidencia-transportada-minima}.

\begin{theorem}[Componentes variacionales de la extensión mínima]
\label{vii:vii:thm:componentes-minimo-conexion}
Para la acción de pantalla
\(S_{\min}=\mathfrak a_m E\), donde \(\mathfrak a_m\) es su
normalización en la recta de acción, escribamos
\[
 b_j=\partial_jb,\qquad \chi_j=\partial_j\chi,
 \qquad a_j=\partial_j\mathfrak a_m.
\]
Con la convención \(\delta S_{\min}=-\tfrac12\sum_j s_j\delta\omega_j\),
las componentes son
\begin{equation}
 \begin{split}
 s_j={}&4\mathfrak a_m\chi\operatorname{Re}
             \langle y,\mathsf C_jf_{\min}-b_j\rangle\\
      &-2\mathfrak a_m\chi_j\langle b,y\rangle-2a_jE.
 \end{split}
 \label{vii:vii:eq:componentes-minimo-conexion}
\end{equation}
Para la normalización cúbica, el segundo término equivale a
\(4\mathfrak a_m E\,\partial_j\ell_m/\ell_m\).
Si el emparejamiento celular de esta misma acción tiene matriz
invertible \(\mathsf M_{\rm cel}\), sus coordenadas de corriente
son \(\sigma=\mathsf M_{\rm cel}^{-1}s\).
\end{theorem}

\begin{proof}
La regla del producto da
\(\partial_jS_{\min}=a_jE+\mathfrak a_m\partial_jE\).
Sustituir \eqref{vii:vii:eq:diferencial-minimo-completo} y multiplicar
por \(-2\) prueba \eqref{vii:vii:eq:componentes-minimo-conexion}.
La derivada de \(\chi_m\) produce la expresión de escala.
Finalmente, comparar
\(\delta S_{\min}=-\tfrac12\delta\omega^{\mathsf T}s\)
con \(-\tfrac12\delta\omega^{\mathsf T}\mathsf M_{\rm cel}\sigma\)
para toda variación da \(\mathsf M_{\rm cel}\sigma=s\),
y su inversa determina las coordenadas.
\end{proof}

La construcción recupera un campo y un peso seleccionados conjuntamente
por la incidencia, y evalúa su respuesta a las variaciones de la misma
realización. El dato de frontera \(b\), la familia de rutas, la escala
\(\ell_m\), la normalización \(\mathfrak a_m\) y el emparejamiento
celular conservan sus reglas de publicación. La interpretación como
corriente material utiliza esos objetos y la acción completa. La
respuesta de Cartan y su densidad métrica posterior se componen con
ella en el dominio de materia correspondiente, conservando cualquier
dependencia de la corriente respecto de la conexión.
